% !PRINTED 376
\weberpage{376}
%<<CONT
\weq{10}{S_n \geqq S' \geqq 0,}
worin aber das Gleichheitszeichen in der unteren Grenze nur für $m = 1$ und in der oberen Grenze nur für $m = 2$ und $b_1 = 1$, also überhaupt nur, wenn $S_n = 1$ ist, vorkommen kann, und dies kann man nach (7) vermeiden, wenn man $n$ gross genug annimmt.

Da aber nach \S.\,118, III durch die Bedingungen (9), (10) die Zahlen $S', R'$ völlig bestimmt sind, und $S_{n-1}, R_{n-1}$ nach (6), (7) denselben Bedingungen genügen, so folgt
\[
S' = S_{n-1}, \qquad R' = R_{n-1}.
\]

Hieraus ergiebt sich nun weiter, dass der Kettenbruch mit der Schlusszahl $x_n$
\weq{11}{(b_0, b_1, b_2 \ldots b_{m-1}, x_n)}
den Werth hat
\[
\dfrac{R_n x_n + R_{n-1}}{S_n x_n + S_{n-1}},
\]
also mit $y$ übereinstimmt.

Wenn man nun $x_n$ weiter in einen Kettenbruch entwickelt
\[
x_n = (a_n, a_{n+1}, a_{n+2} \ldots),
\]
so erhält man aus (2) und (11):
\[
x = (a_0, a_1, a_2 \ldots a_{n-1}, a_n, a_{n+1}, a_{n+2} \ldots)
\]
\[
y = (b_0, b_1, b_2 \ldots b_{m-1}, a_n, a_{n+1}, a_{n+2} \ldots),
\]
oder in Worten ausgesprochen den Satz:

\gesperrt{VI. Die Kettenbruchentwickelungen zweier äquivalenter Zahlen stimmen von einem gewissen Theilnenner an mit einander überein.}

Wir können noch hinzufügen, dass, wenn die Aequivalenz eine eigentliche ist ($\varepsilon = +1$), die Zahl der den übereinstimmenden vorangehenden Theilnenner in beiden eine gerade, oder in beiden eine ungerade ist, und dass, wenn die Aequivalenz uneigentlich ist, diese Zahl in der einen eine gerade, in der anderen eine ungerade ist.

Dass der Satz auch umgekehrt gilt, ist leicht einzusehen. Denn wenn zwei Kettenbrüche, von einem gewissen Theilnenner an übereinstimmen, so können sie so geschrieben werden, dass sie dieselbe Schlusszahl haben. Nun ist der Werth eines Kettenbruches aber immer äquivalent mit jeder seiner Schlusszahlen [\S.\,117, (2)]; \gesperrt{also sind auch zwei Kettenbrüche mit gleicher Schlusszahl unter einander äquivalent.}
